\begin{problem}[第30届全国中学生物理竞赛复赛][匀质细杆转动]{3}
一质量为 $m$、长为 $L$ 的匀质细杆，可绕过其一端的光滑水平轴 $O$ 在竖直平面内自由转动。杆在水平状态由静止开始下摆，
\begin{subproblem}
令 $\lambda = \frac{m}{L}$ 表示细杆质量线密度。当杆以角速度 $\omega$ 绕过其一端的光滑水平轴 $O$ 在竖直平面内转动时，其转动动能可表示为
\begin{equation*}
E_{k} = k \lambda^{\alpha} \omega^{\beta} L^{\gamma}
\end{equation*}
式中，$k$ 为待定的没有单位的纯常数。已知在同一单位制下，两物理量当且仅当其数值和单位都相等时才相等。由此求出 $\alpha$、$\beta$ 和 $\gamma$ 的值。
\end{subproblem}
\begin{subproblem}
已知系统的动能等于系统的质量全部集中在质心时随质心一起运动的动能和系统在质心系（随质心平动的参考系）中的动能之和，求常数 $k$ 的值。
\end{subproblem}
\begin{subproblem}
试求当杆摆至与水平方向成 $\theta$ 角时在杆上距 $O$ 点为 $r$ 处的横截面两侧部分的相互作用力。重力加速度大小为 $g$。
\end{subproblem}
\end{problem}